\scanpage{136}
% archive sheet number: 135

$\mathrm{Az}_P=$ structure de $k$-algèbre de H.C.
sur $P=H\otimes H'$, induisant sur $1\otimes H'$
et $H\otimes1$ les structures d'alg. données.

\diagram{\resizebox{\linewidth}{!}{\begin{tikzcd}[ampersand replacement=\&,column sep=large,row sep=huge]
\mathrm{Az}_P\arrow[r,"{\substack{\alpha\mapsto\det_\alpha\\\deleted{\unsure{\alpha\mapsto\det_\alpha}}}}"]
\arrow[drr,"{\alpha\longrightarrow\mathrm{Tr}_\alpha}"',bend right=7]
\& Q_P\arrow[r,"{\Delta_P\mapsto\Delta_P\mid B}"]
\arrow[dr,dashed,"{\Delta_P\mapsto(U\mapsto\psi_{\Delta_P}(1,U))}"]
\arrow[dr,phantom,"{\text{\unsure{comm.}}}"',near start]
\& Q_{\underline B}\arrow[r,"{\Delta_B\to\varphi_{\Delta_B}}"]
\& \Psi_{\underline B}\arrow[r,"{\varphi\mapsto\psi\ \text{t.q.}\ \varphi+\psi=T_B\otimes T_B}"]
\& \Phi_B\arrow[dll,"{\varphi\mapsto(\underset{\substack{\rotatebox{90}{$\in$}\\H}}{x}\otimes\underset{\substack{\rotatebox{90}{$\in$}\\H'}}{y}\mapsto\varphi(x,y\unsure{'|{}}))}"']\\
\& \& L_P \& \&
\end{tikzcd}}}

La commutativité du triangle total
est immédiate, ainsi que celle du
triangle de gauche, on en conclut
celle du triangle de droite (qui
n'était pas évidente sinon).

\deleted{L'on} On aurait pu aussi
introduire $\underline\Psi_P$ et $\underline\Phi_P$ (en les définissant
avec soin) mais passons ! Parmi les
6 \deleted{\unsure{bases}} ens, celui qui est de
description alg. la plus simple est
sans doute $L_P\subset\mathrm{Bil}(H,H';k)$. On
appellera un $T_P\in L_P$ un \uline{twist}\deleted{\unsure{ing}}
des algèbres \deleted{\unsure{d'Az}} quadratiques
$H,H'$.

% slanted
\margin{C'est défini en termes de $H,H'$ munis \unsure{seulemt|seulement} de $1_H,1_{H'}$ et de $T_H,T_{H'}$, à l'exclusion de la donnée de $\Delta_H,\Delta_{H'}$}

L'objet le \uline{plus fort}
qui en résulte, est la structure
d'alg. de H.C. sur $P=H\otimes H'$,

